\begin{definition}[The structure conditions (E) and (S)]
Let $m\in\mathbb N$, let $A:\mathbb R^m\to\mathbb R^{m\times m}$ and
$H:\mathbb R^m\to\mathbb R$. We say that $(A,H)$ \emph{satisfies the structure conditions},
written $\mathrm{StructureConditionsES}\ m\ A\ H$, if
\begin{enumerate}
\item for every $\xi\neq0$ the matrix $A(\xi)$ is symmetric, $A(\xi)^{T}=A(\xi)$;
\item for every $\xi\neq0$ the matrix $A(\xi)$ is positive semidefinite:
  $\langle A(\xi)\eta,\eta\rangle=\sum_{i}\sum_{j}A_{i,j}(\xi)\eta_j\eta_i\ge0$ for all
  $\eta\in\mathbb R^m$;
\item $A$ is continuous on $\mathbb R^m\setminus\{0\}$;
\item $H$ is continuous on $\mathbb R^m\setminus\{0\}$;
\item (E), (1.9): $\mathcal E(\xi):=\langle A(\xi)\xi,\xi\rangle>0$ for every
  $\xi\in\mathbb R^m\setminus\{0\}$;
\item (S), (1.10): there is a function
  $\kappa:[1,\infty)\times(\mathbb R^m\setminus\{0\})\to(0,\infty)$ such that
  \[
   A(t\xi)=\kappa(t,\xi)\,A(\xi)\quad\text{and}\quad H(t\xi)=t\,\kappa(t,\xi)\,H(\xi)
   \qquad\text{for all }t\ge1\text{ and }\xi\in\mathbb R^m\setminus\{0\}.
  \]
\end{enumerate}

These are exactly the hypotheses imposed on $A$ and $H$ in Section Section~1.2 for the
equation (1.8), $\mathcal Lu=-\mathrm{Tr}(A(\mathcal Xu)\mathcal X\mathcal X^{s}u)
+H(\mathcal Xu)=0$: items (1)--(4) are the standing requirements
$A_{i,j},H\in C(\mathbb R^m\setminus\{0\})$ with $A(\xi)$ symmetric and positive
semidefinite; (5) is the strict ellipticity in the gradient direction (E); and (6) is the
homogeneity along rays (S), which is the form taken, for the operators (1.8), by the
scaling hypothesis of the strong maximum principle of Bardi and Goffi. Nothing is required
at $\xi=0$, and no local Lipschitz continuity is required: the viscosity sub- and
super-solutions of (1.8) are defined with test functions whose horizontal gradient does
not vanish at the point of contact (Definition Definition~2.1), so $A$ and $H$ are never
evaluated at the origin, and the exactness of the relabelling identity under (S) is what
replaces the Lipschitz hypothesis (BB1) of Barles--Busca. In (6), $H$ is required to scale
exactly, with the factor $t\,\kappa(t,\xi)$ and not merely an inequality, because both the
sub-solution and the super-solution side of the scaling hypothesis are used (Remark
Remark~1.1(b)). The positivity of $\kappa$ is part of (S); note that (E) determines it,
$\kappa(t,\xi)=\mathcal E(t\xi)/(t^{2}\mathcal E(\xi))$, so no generality is gained or lost
by quantifying over $\kappa$ existentially.
\end{definition}
